\section{小規模企業共済}

\begin{fillblank}[年金1 2022 問題1(6)(イ)改題]
  小規模企業共済制度の共済金の額に関する説明は以下のとおり。
  \begin{itemize}
    \item 共済金の額は、 （　　A　　） と （　　B　　） の合計金額となる。
    \item （　　A　　） とは、掛金月額、掛金納付月数に応じて、共済事由ごとに小規模企業共済法施
行令で定められている金額である。
\item （　　A　　） の予定利率は、2004 年（平成16 年）4 月以降、それまでの2.5％から
（　　C　　） ％に変更されている。
\item （　　B　　） とは、毎年度の運用収入等に応じて、経済産業大臣が毎年度定める率により算
定される金額である。
\item  （　　B　　） の \warekiyear 年度の支給率は （　　D　　）％ となっている。
  \end{itemize}
\end{fillblank}

\begin{answerbox}
  \begin{enumerate}[A:]
    \item 基本共済金
    \item 付加共済金
    \item 1.0
    \item 0.711（毎年変わる）
  \end{enumerate}
\end{answerbox}

\begin{explanationbox}
  令和8年度付加共済金の支給率について

  \url{https://www.smrj.go.jp/kyosai/skyosai/news/2025/hkj3i80000003n6o.html}
\end{explanationbox}

\begin{shortanswer}[年金1 2021 問題2(4)]
  小規模企業共済制度の設計および税制について、次の(1)および(2)の設問に答えなさい。
  \begin{enumerate}[(1)]
    \item 共済金Ａまたは共済金Ｂについて、その全部または一部が分割払いで支給されるための要件
を簡記しなさい。
\item  解約手当金の税法上の取り扱いについて簡記しなさい。
  \end{enumerate}

\end{shortanswer}

\begin{answerbox}
  \begin{enumerate}[(1)]
    \item 以下の条件を全て満たす場合、共済契約者の請求により、分割払が可能。
    \begin{itemize}
      \item 共済金額が300 万円以上（分割払いと一時払いの併用の場合は330 万円以上）
      \item 請求事由が生じた時点で満60 歳以上
      \item 分割払いと一時払いの併用の場合は、分割払いの合計額が 300 万円以上で、かつ一時払いの額
      が30 万円以上
    \end{itemize}
    \item 一時所得として課税される。ただし、65 歳以上の任意解約等、退職所得となる場合もある。

  \end{enumerate}
\end{answerbox}

\begin{shortanswer}[年金1 2025 問題2(5)]
  小規模企業共済制度の共済金（一時払い） 、共済金（分割払い）、共済金（死亡）、準共済金およ
び解約手当金の税法上の取扱い（所得税の場合は所得の種類とし、所得税以外の場合は税の種
目とする。）について簡潔に入力しなさい。 （150 字以内）
\end{shortanswer}

\begin{answerbox}
  \begin{itemize}
    \item 共済金（一時払い）：退職所得
\item 共済金（分割払い）：雑所得（公的年金等控除の対象）
\item 共済金（死亡）：相続税
\item 準共済金：退職所得
\item 解約手当金：一時所得（ただし、６５歳以上の任意解約は退職所得）
  \end{itemize}
\end{answerbox}

\newpage